\documentclass{ximera}

\title{Algebraic End Behavior}
\author{Kenneth Berglund}

\begin{document}
\begin{abstract}
    Analyzing end behavior of a function algebraically.
\end{abstract}
\maketitle

Consider the polynomial \(g\left(x\right)=2x^5-5x^4-30x^3+5x^2+88x+60\).

\begin{enumerate}
    \item Identify the degree of the polynomial.
    \item Which of the 6 terms, \(2x^5\), \(5x^4\), \(30x^3\), \(5x^2\),
    \(88x\), or \(60\), is greatest when:
    \begin{enumerate}
        \item \(x = 0\)
        \item \(x = 1\)
        \item \(x = 3\)
        \item \(x = 5\)
    \end{enumerate}
    \item Write a sentence describing what happens to \(g\left(x\right)\) 
    as \(x\) gets bigger and bigger in the positive direction. 
    Remember that \(g(x)\) is the \(y\)-value associated with \(x\). 
    An example sentence might be ``As \(x\) gets bigger and bigger in 
    the positive direction, \(g(x)\) gets...''.
    \item Write a sentence describing what happens to \(g\left(x\right)\) 
    as \(x\) gets bigger and bigger in the negative direction. 
\end{enumerate}


\end{document}
